%%%
%%% Radical "⽏"
%%%
\section*{Radical 80: ``⽏'' (母)}\addcontentsline{toc}{section}{Radical 80: ⽏,母}\addcontentsline{loh}{figure}{\#\#\#\# 80: ⽏}

%%%%%%%%%% 母 %%%%%%%%%%
\subsection*{母}\addcontentsline{loh}{figure}{母}

\begin{Entry}{母}{5}{⽏}[Kangxi 80]
  \begin{Phonetics}{母}{mu3}[][HSK 6]
    \definition*{s.}{Sobrenome: Mu}
    \definition{adj.}{fêmea}
    \definition[位,名,个,些]{s.}{mãe | fêmea (animal) | origem; pais | parentes idosas; geralmente se refere a mulheres idosas | côncavo | fonte; algo que tem a capacidade ou função de produzir outras coisas}
  \antonymref{公}{gong1}
  \end{Phonetics}
\end{Entry}

\begin{Entry}{母女}{5,3}{⽏,⼥}
  \begin{Phonetics}{母女}{mu3-nv3}[][HSK 6]
    \definition{s.}{mãe e filha}
  \end{Phonetics}
\end{Entry}

\begin{Entry}{母子}{5,3}{⽏,⼦}
  \begin{Phonetics}{母子}{mu3zi3}[][HSK 6]
    \definition{s.}{mãe e filho}
  \end{Phonetics}
\end{Entry}

\begin{Entry}{母鸡}{5,7}{⽏,⿃}
  \begin{Phonetics}{母鸡}{mu3ji1}[][HSK 6]
    \definition{s.}{galinha}
  \end{Phonetics}
\end{Entry}

\begin{Entry}{母亲}{5,9}{⽏,⼇}
  \begin{Phonetics}{母亲}{mu3qin5}[][HSK 3]
    \definition[位,名,个,些]{s.}{mãe}
  \end{Phonetics}
\end{Entry}

\begin{Entry}{母语}{5,9}{⽏,⾔}
  \begin{Phonetics}{母语}{mu3yu3}
    \definition{s.}{língua materna; a primeira língua que uma pessoa aprende, geralmente a língua padrão de seu povo nativo ou um dialeto específico | língua"-mãe; a origem comum de determinadas línguas}
  \synonymref{方言}{fang1yan2}
  \antonymref{外文}{wai4wen2}
  \antonymref{外语}{wai4yu3}
  \end{Phonetics}
\end{Entry}

%%%%%%%%%% 每 %%%%%%%%%%
\subsection*{每}\addcontentsline{loh}{figure}{每}

\begin{Entry}{每}{7}{⽏}
  \begin{Phonetics}{每}{mei3}[][HSK 3]
    \definition{adv.}{cada um; cada qual; indica qualquer uma das repetições ou um conjunto de repetições de um movimento}
    \definition{pron.}{cada; cada um; cada qual; refere"-se a qualquer indivíduo do grupo, enfatizando as semelhanças entre os indivíduos}
  \end{Phonetics}
\end{Entry}

\begin{Entry}{每个}{7,3}{⽏,⼈}
  \begin{Phonetics}{每个}{mei3ge4}
    \definition{pron.}{cada; cada um}
  \end{Phonetics}
\end{Entry}

\begin{Entry}{每个人}{7,3,2}{⽏,⼈,⼈}
  \begin{Phonetics}{每个人}{mei3ge5 ren2}
    \definition{s.}{todo mundo; qualquer pessoa}
  \end{Phonetics}
\end{Entry}

\begin{Entry}{每天}{7,4}{⽏,⼤}
  \begin{Phonetics}{每天}{mei3tian1}
    \definition{s.}{todos os dias; cada dia (único)}
  \synonymref{经常}{jing1chang2}
  \synonymref{时常}{shi2chang2}
  \synonymref{天天}{tian1tian1}
  \end{Phonetics}
\end{Entry}

\begin{Entry}{每当}{7,6}{⽏,⼹}
  \begin{Phonetics}{每当}{mei3dang1}[][HSK 7-9]
    \definition{adv.}{sempre que; todas as vezes; toda vez que isso acontece; em qualquer momento}
  \end{Phonetics}
\end{Entry}

\begin{Entry}{每次}{7,6}{⽏,⽋}
  \begin{Phonetics}{每次}{mei3ci4}
    \definition{adv.}{cada vez; todas as vezes; sempre}
  \synonymref{常常}{chang2chang2}
  \synonymref{老是}{lao3shi4}
  \synonymref{每逢}{mei3feng2}
  \synonymref{时常}{shi2chang2}
  \antonymref{从不}{cong2bu4}
  \antonymref{偶尔}{ou3'er3}
  \antonymref{偶然}{ou3ran2}
  \end{Phonetics}
\end{Entry}

\begin{Entry}{每逢}{7,10}{⽏,⾡}
  \begin{Phonetics}{每逢}{mei3feng2}[][HSK 7-9]
    \definition{adv.}{sempre que; em todas as ocasiões; em todas as situações; toda vez que me deparo com isso; toda vez que chego}
  \end{Phonetics}
\end{Entry}

%%%%%%%%%% 毒 %%%%%%%%%%
\subsection*{毒}\addcontentsline{loh}{figure}{毒}

\begin{Entry}{毒}{9}{⽏}
  \begin{Phonetics}{毒}{du2}[][HSK 5]
    \definition*{s.}{Sobrenome: Du}
    \definition{adj.}{veneno; toxina; propriedade ou substância prejudicial aos organismos vivos | droga; narcóticos | vírus; vírus de computador | influência venenosa}
    \definition{adj.}{venenoso; tóxico; envenenado | malicioso; cruel; feroz}
    \definition{v.}{matar com veneno; envenenar | envenenar (a mente de alguém)}
  \end{Phonetics}
\end{Entry}

\begin{Entry}{毒杀}{9,6}{⽏,⽊}
  \begin{Phonetics}{毒杀}{du2sha1}
    \definition{v.}{matar com veneno; envenenar; matar por envenenamento | eliminar os percevejos com substâncias tóxicas}
  \end{Phonetics}
\end{Entry}

\begin{Entry}{毒物}{9,8}{⽏,⽜}
  \begin{Phonetics}{毒物}{du2wu4}
    \definition{s.}{substância venenosa; veneno | substância tóxica | toxina | peçonha}
  \end{Phonetics}
\end{Entry}

\begin{Entry}{毒品}{9,9}{⽏,⼝}
  \begin{Phonetics}{毒品}{du2pin3}[][HSK 6]
    \definition[种,点]{s.}{drogas; veneno; narcóticos; refere"-se ao ópio, morfina, heroína, etc. usados como vício}
  \end{Phonetics}
\end{Entry}

\begin{Entry}{毒害}{9,10}{⽏,⼧}
  \begin{Phonetics}{毒害}{du2hai4}
    \definition{v.}{envenenar (prejudicar com uma substância tóxica) | envenenar (as mentes das pessoas)}
    \definition{v.}{matar por envenenamento; envenenar | envenenar (a mente de alguém); macular; causar dano | infectar}
  \synonymref{伤害}{shang1hai4}
  \synonymref{损害}{sun3hai4}
  \synonymref{中毒}{zhong4du2}
  \antonymref{保护}{bao3hu4}
  \antonymref{解救}{jie3jiu4}
  \antonymref{救助}{jiu4zhu4}
  \end{Phonetics}
\end{Entry}

\begin{Entry}{毒蛇}{9,11}{⽏,⾍}
  \begin{Phonetics}{毒蛇}{du2she2}
    \definition{s.}{cobra venenosa (ou peçonhenta); víbora; serpente; toxicofídia; serpentes venenosas geralmente têm cabeças triangulares e glândulas que secretam veneno; quando uma serpente venenosa morde uma pessoa ou animal, o veneno flui de suas presas, causando envenenamento}
  \end{Phonetics}
\end{Entry}

%%%%% EOF %%%%%

